% !TeX program = xelatex
% Build: xelatex -> bibtex -> xelatex -> xelatex
\documentclass[runningheads,orivec]{llncs}

\usepackage{amsmath,amssymb}
\usepackage{booktabs}
\usepackage{enumitem}
\usepackage{graphicx}
\usepackage{hyperref}
\usepackage{tabularx}
\usepackage{array}
\usepackage{placeins}

\hypersetup{hidelinks}
\setlength{\emergencystretch}{1.5em}
\newcommand{\xwing}{X-Wing}

\title{X-Wing on Cortex-M85:\
Architecture-Aware Integration and End-to-End Evaluation}
\titlerunning{X-Wing on Cortex-M85}

\author{No Author Given}
\authorrunning{No Author Given}
\institute{No Institute Given}

\begin{document}
\raggedbottom
\maketitle

\begin{abstract}
\xwing{} combines ML-KEM-768, a post-quantum key encapsulation mechanism (KEM),
with X25519, a classical elliptic-curve key-agreement scheme. Optimized
implementations exist for both components, but improving individual functions
does not necessarily improve the complete \xwing{} API on a microcontroller
with limited issue width and register resources. This paper preserves the
cryptographic specification and API of \xwing{} draft-10 while jointly adapting
the bottlenecks of both computational regimes to the execution resources of the
Cortex-M85.

The final integrated implementation comprises eight optimizations. For ML-KEM,
we use the 128-bit M-profile Vector Extension (MVE) for the Number Theoretic
Transform (NTT), match the NTT storage order to downstream operations, and
vectorize 11 data-format transformations. For X25519, we implement a width-3
signed comb for basepoint 9 with a 33,024-byte precomputation table, reschedule
field multiplication and squaring, and implement the conditional swap with MVE.
We also replace two field inversions that coexist during \xwing{} encapsulation
with one zero-safe batch inversion and exploit alignment information in common
data-copy and initialization routines.

Each candidate was evaluated first in isolation and then in the actual \xwing{}
API before integration. The reference implementation and the final integrated
implementation were alternated within one executable. We performed five
independent experiments, reflashing the same board image before each experiment.
Across ten paired comparisons, the minimum observed cycle reductions were
19.52\% for key generation, 19.75\% for encapsulation, 11.66\% for warm
decapsulation with reuse of an expanded key, and 15.48\% for cold decapsulation
including key expansion from a 32-byte seed. Both implementations passed
standard test vectors, byte-for-byte equivalence checks, invalid-ciphertext and
low-order-point tests, and stack-boundary checks. The contribution is not a new
cryptographic primitive. It is the selection of where existing techniques apply
within \xwing{}, their Cortex-M85-specific implementation, and an end-to-end
demonstration that function-level gains persist at the complete API level.

\keywords{\xwing{} \and hybrid KEM \and ML-KEM \and X25519 \and Cortex-M85
\and M-profile Vector Extension \and embedded cryptographic optimization}
\end{abstract}

\section{Introduction}\label{sec:introduction}

During the transition to post-quantum cryptography, hybrid key establishment is
important because it combines classical and post-quantum cryptography instead
of relying exclusively on one new algorithm. \xwing{} is a hybrid KEM that
combines the two shared secrets obtained from X25519 and ML-KEM-768 with
SHA3-256 \cite{XWing24,XWingDraft10}. The construction is designed so that one
component continues to provide security if the other becomes vulnerable.

On embedded devices, however, optimizations of ML-KEM and X25519 must be
evaluated separately to determine whether they translate into performance gains
for the complete \xwing{} API. We target the Cortex-M85 to evaluate this
translation on a high-performance M-profile microcontroller equipped with the
Arm M-profile Vector Extension (MVE, also known as Helium). ML-KEM exposes
substantial data parallelism that can use 128-bit MVE, whereas X25519 field
arithmetic contains long scalar dependency chains. The platform is therefore
well suited to studying the integration effects that arise when vector-oriented
and scalar-oriented computations share the same execution resources. Limited
instruction issue width, a small scalar register file, and finite load/store
bandwidth nevertheless prevent the assumption that a faster function will
necessarily accelerate the complete \xwing{} implementation.

\subsection{Related Work}

The original \xwing{} work presented the construction and its security rationale
and evaluated the performance of optimized implementations \cite{XWing24}. To
retain implementation simplicity, that work treated both components as black
boxes, while also discussing further optimization opportunities across internal
ML-KEM processing and the \xwing{} key-derivation stage. This paper targets the
subsequently specified draft-10 algorithm and API \cite{XWingDraft10}.

For ML-KEM implementations, pqm4 provides optimized Cortex-M4 implementations
and a common evaluation framework \cite{pqm4}. Prior work reduced the cost of
major operations through assembly NTTs, modular arithmetic, and lower memory
usage. Work on Helium NTTs presented MVE implementations of the NTT and
polynomial multiplication on Cortex-M55 \cite{HeliumNTT22}. pqmx extends this
line of MVE implementations \cite{pqmx}, while SLOTHY uses constraint solving
to reschedule cryptographic assembly for a target microarchitecture
\cite{Slothy22,SlothyM7_25}.

For X25519, prior studies have developed constant-time Montgomery ladders and
optimized field multiplication and squaring for small microcontrollers and the
Cortex-M family \cite{RFC7748,Lenngren,FujiiAranha17,DullEtAl15}. A comb using a
precomputation table when the input point is fixed to basepoint 9 and a batch
inversion that replaces several denominator inversions with one inversion are
also established elliptic-curve implementation techniques \cite{GuideECC04}.

This work does not introduce these algorithms, but it also goes beyond simply
linking existing code. We connect an MVE NTT to the coefficient representation
and downstream storage order of the existing ML-KEM implementation, vectorize
11 data-format transformations and the X25519 conditional swap, implement and
reschedule a fixed-base comb and field arithmetic for the Cortex-M85, and apply
a zero-safe batch inversion where two X25519 outputs coexist during \xwing{}
encapsulation. The contribution boundary is the \xwing{}-specific placement of
these techniques, their Cortex-M85-specific implementation, and end-to-end
validation of the complete API on the same board.

\subsection{Contributions}

This paper makes the following contributions.
\begin{enumerate}[leftmargin=*,itemsep=2pt]
\item We connect an MVE NTT from the pqmx implementation lineage to the
      coefficient representation and downstream storage order of the existing
      ML-KEM implementation, eliminate a separate scalar reordering pass, and
      implement 11 ML-KEM data-format transformations with MVE instructions.
\item We implement a constant-time width-$w=3$ signed comb for the scalar
      multiplication whose input is basepoint 9 in \xwing{}. We also restructure
      the stack usage and instruction schedules of X25519 field multiplication
      and squaring for Cortex-M85, and optimize copies and initialization of
      internal buffers whose sizes and alignments are known.
\item We replace the scalar xor-mask conditional swap with four groups of
      128-bit MVE operations. During encapsulation, we combine the inversions of
      two simultaneously live X25519 outputs and preserve the original outputs
      for zero denominators and low-order input points.
\item In addition to per-candidate evaluation, we place the reference
      implementation with all optimizations disabled and the final integrated
      implementation with all eight optimizations enabled in one executable and
      directly compare them under the same board and measurement procedure.
\end{enumerate}

The direct comparison reduced cycles by at least 19.52\% for key generation,
19.75\% for encapsulation, 11.66\% for warm decapsulation with expanded-key
reuse, and 15.48\% for cold decapsulation including expansion from a 32-byte
seed.

\section{Background and Problem Definition}\label{sec:background}

\subsection{Execution Structure of \xwing{}}

\xwing{} executes ML-KEM-768 and X25519 independently and supplies both results
and context information to a SHA3-256 combiner
\cite{XWingDraft10,FIPS203,RFC7748}.

\paragraph{Key generation.}
The implementation generates an ML-KEM key pair, multiplies the fixed basepoint
9 by an X25519 secret scalar to derive a public key, and assembles the component
keys in the \xwing{} format. We refer to an operation whose input point is
basepoint 9 as fixed-base scalar multiplication.

\paragraph{Encapsulation.}
ML-KEM produces a ciphertext and a shared secret. One ephemeral X25519 secret
scalar is multiplied by basepoint 9 to produce an ephemeral public key and by
the recipient's public key to produce an X25519 shared secret. We refer to the
latter operation, whose input point varies with each recipient public key, as
variable-base scalar multiplication. The results of both algorithms and the
context information are then processed by the SHA3-256 combiner.

\paragraph{Decapsulation.}
The implementation decapsulates the ML-KEM ciphertext and selects the implicit
rejection value when the ciphertext is invalid. It then computes the X25519
shared secret through variable-base scalar multiplication using the sender's
public key and processes both results and the context information with the same
SHA3-256 combiner. We call the condition that reuses an already expanded
2,464-byte decapsulation key \emph{warm decapsulation}. We call the condition
that includes expansion of the decapsulation key from a 32-byte seed
\emph{cold decapsulation}. The latter key expansion regenerates the X25519
public key and therefore includes fixed-base scalar multiplication.

Consequently, the fixed-base optimization applies to key generation,
encapsulation, and cold decapsulation, but not to warm decapsulation. During
encapsulation, the fixed-base ephemeral public key and the variable-base shared
secret are computed concurrently, so the inversions required by the two outputs
can share one batch inversion. We refer to this joint inversion of the two
outputs as \emph{paired inversion}.

\subsection{ML-KEM Bottlenecks}

ML-KEM uses a forward NTT, coefficient-wise operations, and an inverse NTT to
accelerate polynomial multiplication \cite{FIPS203}. It also repeatedly performs
centered binomial distribution (CBD) sampling, serialization, compression and
decompression, message conversion, and comparison. MVE arithmetic can
accelerate the NTT, and independent coefficients can be processed concurrently
in the data-format transformations. If the NTT output order does not match the
input order required by downstream operations, however, an additional scalar
reordering pass is required.

\subsection{X25519 Bottlenecks}

The X25519 Montgomery ladder repeatedly performs field addition, subtraction,
multiplication, squaring, and a conditional swap \cite{RFC7748}. Intermediate
points remain in projective $X/Z$ coordinates, and conversion to the final
affine coordinate requires the inverse of $Z$. Field inversion is much more
expensive than multiplication, so multiple outputs can share one inversion of
the product of their denominators. A precomputation table for basepoint 9 can be
used for fixed-base scalar multiplication, whereas the same table cannot be
used for variable-base scalar multiplication whose input is another party's
public key.

\subsection{Relevant Cortex-M85 Resources}

Cortex-M85 is an in-order Armv8.1-M core that provides 128-bit vector operations
through MVE \cite{ArmM85Opt}. MVE is advantageous when the same operation is
applied to several coefficients or 32-bit words. The following constraints can
nevertheless limit integrated performance.
\begin{itemize}[leftmargin=*,itemsep=1pt]
\item The small scalar register file makes it difficult to keep many
      intermediate values live simultaneously.
\item Memory accesses from MVE and scalar instructions cannot be issued
      concurrently without limit.
\item Instruction dependencies and dual-issue rules can give different cycle
      counts to sequences containing the same number of instructions.
\item If register pressure introduces spills or instructions that regenerate
      constants, their cost can exceed the benefit obtained from otherwise idle
      issue slots.
\end{itemize}

\FloatBarrier
\section{Design}\label{sec:design}

\subsection{Overview}

Table~\ref{tab:transform-scope} summarizes the final integrated implementation,
which enables all eight optimizations. None changes the cryptographic
specification or the external API.

\begin{table}[t]
\centering
\caption{Eight optimizations applied to \xwing{} on Cortex-M85 and the
operations to which they apply.}
\label{tab:transform-scope}
\footnotesize
\renewcommand{\arraystretch}{1.12}
\begin{tabularx}{\textwidth}{@{}cl>{\raggedright\arraybackslash}Xl@{}}
\toprule
No. & Component & Optimization & Applicable operations \\
\midrule
1 & ML-KEM & MVE NTT with downstream-compatible storage order & All \\
2 & ML-KEM & MVE implementation of 11 data-format transformations & All \\
3 & X25519 & Width-$3$ signed comb for basepoint 9 & Key generation, encapsulation, cold \\
4 & X25519 & Fixed stack frame and rescheduling for field multiplication & All X25519 \\
5 & X25519 & Rescheduling for field squaring & All X25519 \\
6 & Common & Alignment-aware data copy and initialization & All \\
7 & X25519 & MVE xor-mask conditional swap & Encapsulation, warm, cold \\
8 & \xwing{}/X25519 & Zero-safe paired inversion & Encapsulation \\
\bottomrule
\end{tabularx}
\end{table}

\begin{figure}[t]
\centering
\includegraphics[width=\textwidth]{assets/fig01_xwing_optimization_map.pdf}
\caption{\xwing{} data flow and scope of the eight retained optimizations.
Numbered badges refer to Table~\ref{tab:transform-scope}. The fixed-base comb is
absent from warm decapsulation, MVE cswap is confined to variable-base ladders,
and zero-safe paired inversion links the two X25519 outputs only during
encapsulation. Alignment-aware copy and zeroing span all four API paths.}
\label{fig:optimization-map}
\end{figure}

\subsection{ML-KEM: NTT Storage Order and Data-Format Transformations}

After storing an NTT result, the reference ML-KEM implementation calls the
scalar function \texttt{rev4} to permute groups of four coefficients into the
order required by downstream operations. It also performs CBD sampling,
serialization, compression, message conversion, and comparison with scalar C
loops. These reordering and repeated coefficient transformations incur
additional memory accesses and scalar instructions beyond the NTT arithmetic.

We connect the forward and inverse Cortex-M85 NTTs from the pqmx implementation
lineage to the Plantard coefficient representation and admissible ranges used
by the reference implementation \cite{pqmx}. We then use the MVE structured
store that pqmx calls \texttt{current-order}, making the forward NTT store
coefficients directly in the order required by downstream base multiplication
or matrix accumulation. This removes the separate scalar reordering pass
between the NTT and downstream operations. We also replace 11 data-format
transformations with fixed-trip-count implementations based on MVE compiler
intrinsics, extracting, quantizing, or comparing several coefficients
concurrently in 128-bit vectors.

These two optimizations apply wherever ML-KEM executes: key generation,
encapsulation, warm decapsulation, and cold decapsulation. Predicate management
for active vector lanes and byte insertion with high data-rearrangement cost
remain scalar. The optimization does not change the NTT or ML-KEM mathematics.
It adapts the coefficient representation, range, and storage order of an
existing MVE NTT to the \xwing{} implementation and vectorizes only the
data-format transformations whose gains remain at the complete API level.

\subsection{X25519: Fixed-Base Scalar Multiplication and Field Arithmetic}

The reference implementation uses the general Montgomery ladder even when the
input point for public-key generation is fixed to basepoint 9. Within the
ladder, \texttt{fe25519\_mul} uses a push/pop frame that saves and restores
registers at function entry and exit and introduces stack-pointer (SP) update
dependencies. Some multiplication and squaring instructions also wait for
earlier results. Repeated work in fixed-base scalar multiplication and
dependency stalls in field arithmetic therefore remain distinct bottlenecks.

For fixed-base scalar multiplication, we use a width-$w=3$ signed comb and a
33,024-byte precomputation table. Instead of secret-dependent branches or a
conditional sign change of a field element, the selected sign is handled by a
constant-time swap of two pointers. Intermediate values in repeated squaring
chains remain in registers where possible. For field multiplication, we flatten
the stack frame into fixed SP offsets and place independent instructions in
dependency stalls. For field squaring, we retain the 94-instruction count while
placing independent operations between instructions on long dependency chains.

The fixed-base comb applies only to key generation, encapsulation, and cold key
expansion; it does not apply to variable-base scalar multiplication or warm
decapsulation, which reuses an already expanded key. In contrast, rescheduling
field multiplication and squaring applies to every operation that executes
X25519. Signed combs and instruction rescheduling are established techniques.
Our contribution is their implementation at the exact \xwing{} call sites and
for the Cortex-M85 execution characteristics, together with an end-to-end
measurement that includes the 33,024-byte flash cost.

\subsection{X25519: MVE Conditional Swap and Paired Inversion}

The reference conditional swap in the variable-base Montgomery ladder processes
16 32-bit words in a scalar xor-mask loop. During \xwing{} encapsulation, one
ephemeral secret scalar produces both a fixed-base ephemeral public key and a
variable-base shared secret, yet the reference implementation separately
inverts the denominators of the two projective-coordinate outputs. The former
operation repeats the same word-level computation, while the latter duplicates
an expensive field inversion.

We replace the conditional swap with four groups of 128-bit MVE load, xor, and,
xor, and store operations. The instruction flow is identical whether the mask
is zero or has all bits set, and all vector temporaries are consumed within the
function. This MVE conditional swap (MVE cswap) applies to encapsulation, warm
decapsulation, and cold decapsulation, where the variable-base ladder executes.

Let $d_1$ and $d_2$ denote the two denominators in encapsulation. For each zero
denominator, we create $\widetilde d_i$ by replacing the denominator with one
only during the computation. The following construction reduces the two field
inversions to one:
\begin{equation}
\begin{aligned}
\widetilde d_i & =
\begin{cases}
1, & d_i=0,\\
d_i, & d_i\neq0,
\end{cases}
\quad i\in\{1,2\},\\
t & = (\widetilde d_1\widetilde d_2)^{-1}, &
u_1 & = \widetilde d_2t, &
u_2 & = \widetilde d_1t .
\end{aligned}
\label{eq:paired-inversion}
\end{equation}

The value $u_i$ is used for coordinate normalization and equals $d_i^{-1}$
when $d_i\neq0$. Zero testing compares the complete normalized field
representation. A saved mask then restores the original all-zero output for a
low-order input point. The tests and output masking follow the same instruction
flow independently of the input. Batch inversion itself is an established
technique. Our implementation contribution is the zero-safe connection of this
technique at the \xwing{} encapsulation point where two X25519 results coexist.
Paired inversion does not apply to key generation or decapsulation, where only
one output is present.

\FloatBarrier
\begin{figure}[t]
\centering
\includegraphics[width=\textwidth]{assets/fig03_paired_inversion.pdf}
\caption{Zero-safe paired inversion in \xwing{} encapsulation. The reference
implementation independently inverts $d_1$ and $d_2$. The final integrated
implementation records a constant-time zero mask for each denominator,
substitutes $\widetilde d_i=1$ only during calculation when $d_i=0$, computes
$t=(\widetilde d_1\widetilde d_2)^{-1}$, and recovers
$u_1=\widetilde d_2t$ and $u_2=\widetilde d_1t$. The saved masks restore the
original all-zero outputs.}
\label{fig:paired-inversion}
\end{figure}

\subsection{Common Memory Operations and Constant-Time Compatibility}

The general-purpose \texttt{memcpy} and \texttt{memset} routines in newlib-nano
support many lengths and alignments, whereas the sizes and alignments of several
internal \xwing{} buffers are known. We use routines that branch only on address
alignment and length and perform copies and initialization one word at a time.
Copies already inlined by the compiler are outside the replacement scope.

None of the eight optimizations changes the cryptographic specification or the
API. Rescheduling and changes to storage order preserve fixed instruction and
address patterns. The fixed-base table selection, MVE cswap, and zero-safe paired
inversion use masks to avoid secret-dependent control flow. The alignment-aware
memory routines branch only on public address alignment and length. Functional
equivalence tests and execution-time tests with fixed and random ciphertexts
support this design, but they do not constitute a leakage proof for power,
electromagnetic emanations, or every cache state.

\FloatBarrier
\section{Evaluation}\label{sec:evaluation}

\subsection{Reference and Final Integrated Implementations}

The unoptimized reference implementation consists of ML-KEM and Keccak from the
pqm4 implementation lineage, X25519 from the Lenngren implementation lineage,
scalar C data-format transformations, and newlib-nano memory routines, all built
for Cortex-M85. The final integrated implementation enables all eight
optimizations from Sect.~\ref{sec:design}. Both implementations use the same
\xwing{} draft-10 specification, inputs, API, and SHA3-256 combination order.

For the primary performance comparison, the reference implementation and the
final integrated implementation are included in one Executable and Linkable
Format (ELF) file, and a public runtime mode selects one implementation. They
therefore share the same link and execution environment, clock configuration,
measurement call sites, and inputs. Experimental selection and instrumentation
code is common to both executions. We consequently interpret only the
difference measured within the same executable rather than treating the
absolute cycle counts as the performance of a deployable library.

\subsection{Device and Build}

\begin{itemize}[leftmargin=*,itemsep=1pt]
\item Device: EK-RA8M1, Cortex-M85, 480 MHz
\item Compiler: GCC 13.2.1, \texttt{-O2}
\item Code placement: instruction tightly coupled memory (ITCM)
\item Data and stack placement: data tightly coupled memory (DTCM)
\item Cycle counter: Data Watchpoint and Trace CYCCNT (DWT CYCCNT)
\item Each implementation-operation combination: median of 100 measurements
\item Independent repetition: reflash the same board image and measure from boot
\end{itemize}

A write-stall interval observed at one specific ITCM address on the evaluation
board was avoided with link padding. Before every run, we read the flash contents
back and confirmed that they matched the recorded image. Only runs that passed
the known-answer tests (KATs) were included in the results.

\subsection{Measurement Order}

The final comparison executes the two implementations in one executable in the
following order:
\begin{center}
Reference $\rightarrow$ Final integrated $\rightarrow$ Final integrated
$\rightarrow$ Reference
\end{center}
This order counterbalances effects caused by execution order. We performed five
independent repetitions, reflashing the same board image before each repetition.
Within each repetition, adjacent reference and final integrated executions form
two pairs, yielding ten paired comparisons per operation. Let
$C_{\mathrm{ref},i}(o)$ and $C_{\mathrm{final},i}(o)$ denote the cycles of the
reference and final integrated implementations, respectively, for the $i$th
paired comparison of operation $o$. The improvement and reported value are
defined as
\begin{equation}
r_i(o)=100\left(1-\frac{C_{\mathrm{final},i}(o)}{C_{\mathrm{ref},i}(o)}\right),
\qquad
r_{\mathrm{reported}}(o)=\min_{1\le i\le10}r_i(o).
\label{eq:reported-rate}
\end{equation}
We report both the range across the ten improvements and the value defined in
Eq.~\eqref{eq:reported-rate}.

\subsection{Measured Operations}

We measure key generation, encapsulation, warm decapsulation, and cold
decapsulation. Warm decapsulation reuses an expanded 2,464-byte decapsulation
key. Cold decapsulation first expands the key from a compressed 32-byte seed.
Because these conditions execute different work, we do not combine them into a
single average for decapsulation.

\subsection{Correctness Criteria}

Only candidates that passed all of the following checks were included in the
performance results.
\begin{itemize}[leftmargin=*,itemsep=1pt]
\item RFC 7748 X25519 KATs and SHA3-256 KATs
\item ML-KEM-768 encapsulation-decapsulation round trips and implicit rejection
      of invalid ciphertexts
\item \xwing{} draft-10 Appendix C test vectors
\item Byte-for-byte comparison of public keys, secret keys, ciphertexts, and
      shared secrets from both implementations on eight fixed inputs
\item Agreement between warm and cold results and tests with three low-order
      X25519 input points
\item Stack-boundary corruption checks and comparison of the recorded image with
      the actual flash contents
\end{itemize}

\FloatBarrier
\subsection{Candidate Selection and Contribution Analysis}

Each optimization candidate is evaluated at three levels. First, we measure
only the target function to establish feasibility and its local effect. Second,
we connect the candidate to the actual caller and remeasure the complete key
generation, encapsulation, and decapsulation APIs. Finally, we verify its effect
again in a diagnostic configuration containing the retained optimizations and
in the executable that compares the reference and final integrated
implementations.

\begin{figure}[t]
\centering
\includegraphics[width=\textwidth]{assets/fig02_adoption_pipeline.pdf}
\caption{Candidate adoption and reporting procedure. Each change first passes
correctness and function-level measurement, is then integrated into a real
\xwing{} call path, and is remeasured at API level. Retained changes are checked
in a combined diagnostic build. Intermediate and leave-one-out percentages
remain diagnostic; the reported result is the minimum of ten paired
improvements from the final same-executable comparison between the reference
and final integrated implementations.}
\label{fig:adoption-pipeline}
\end{figure}

We quantify the individual contributions of the six-optimization diagnostic by
removing one optimization at a time from the configuration containing all six.
We refer to this procedure as leave-one-out analysis. Intermediate and final
experiments can use different executables and measurement call sites; therefore,
we neither add nor multiply improvements obtained from different experiments.

\FloatBarrier
\section{Results}\label{sec:results}

\subsection{Final Performance}

Table~\ref{tab:final-performance} presents the final end-to-end results.
\begin{table}[t]
\centering
\caption{End-to-end cycles of the reference and final integrated
implementations alternated within one executable. Each cycle count is the
median of 100 measurements. The ranges summarize ten paired comparisons from
five independent reflashes. The reported value is the minimum observed
improvement defined in Eq.~\eqref{eq:reported-rate}.}
\label{tab:final-performance}
\scriptsize
\setlength{\tabcolsep}{2.2pt}
\begin{tabular}{@{}lrrrr@{}}
\toprule
Operation & Reference cycles & Integrated cycles & Improvement range & Reported \\
\midrule
Key generation & 870,869--870,970 & 700,834--700,865 & 19.521--19.534\% & \textbf{19.52\%} \\
Encapsulation & 1,326,924--1,326,931 & 1,064,853--1,064,862 & 19.750--19.751\% & \textbf{19.75\%} \\
Warm decapsulation & 934,803--934,837 & 825,746--825,793 & 11.661--11.670\% & \textbf{11.66\%} \\
Cold decapsulation & 1,802,626--1,802,655 & 1,523,463--1,523,502 & 15.484--15.488\% & \textbf{15.48\%} \\
\bottomrule
\end{tabular}
\end{table}

These results are the final effects of applying all eight optimizations to the
complete \xwing{} API. The subsequent tables are diagnostic evidence explaining
where the final effects arise; they are not quantities to be added to the final
results. All five independent repetitions used the same ELF file and the same
S-record (SREC) board image. We verified the flash contents after every
programming step, and every run passed all functional and stack-boundary checks.
The narrow improvement ranges show that the difference between the reference
and final integrated implementations is reproducible under the same board and
measurement conditions.

\subsection{Integrated Effect}

We first performed a separate diagnostic experiment that combined six of the
final eight optimizations: ML-KEM NTT and storage order, ML-KEM data-format
transformations, the fixed-base comb, field multiplication, field squaring, and
alignment-aware memory routines. This experiment also placed the configuration
with all six disabled and the configuration with all six enabled in one
executable for a direct comparison. Key generation, encapsulation, warm
decapsulation, and cold decapsulation improved by 25.96\%, 20.25\%, 11.22\%, and
18.30\%, respectively.

These values are direct comparisons between configurations before and after all
six optimizations; they are not sums of six function-level improvements. The
six-optimization diagnostic and the final direct comparison use different
executables and common measurement code. We therefore do not add the
improvements of the two additional optimizations to values such as 25.96\% to
construct the final improvement.

\FloatBarrier
\subsection{Individual Contributions}

Table~\ref{tab:leave-one-out} reports the cycle increases observed when each
optimization was removed individually from the six-optimization diagnostic.
\begin{table}[t]
\caption{Leave-one-out results for the six-optimization diagnostic. Each value
is the cycle increase caused by removing one optimization and must not be added
to the improvements from the final direct comparison.}
\label{tab:leave-one-out}
\centering
\footnotesize
\setlength{\tabcolsep}{3.0pt}
\resizebox{\textwidth}{!}{%
\begin{tabular}{@{}lrrrr@{}}
\toprule
Optimization & Key generation & Encapsulation & Warm & Cold \\
\midrule
Fixed-base comb & 152,860 & 152,710 & 16 & 152,732 \\
NTT and storage order & 22,556 & 18,002 & 28,180 & 50,808 \\
MVE data-format transformations & 11,120 & 33,300 & 39,386 & 50,552 \\
Alignment-aware memory routines & 12,660 & 10,042 & 9,906 & 21,380 \\
Field-squaring rescheduling & 3,488 & 22,770 & 19,296 & 22,794 \\
Field-multiplication rescheduling & 4,182 & 11,950 & 7,632 & 11,898 \\
\bottomrule
\end{tabular}
}
\end{table}

The fixed-base comb makes the largest individual contribution to key
generation, encapsulation, and cold decapsulation, while its effect on warm
decapsulation is at the noise level. This result agrees with the applicability
described in Sect.~\ref{sec:background}. The NTT and data-format transformations
contribute to all four measured operations in which ML-KEM executes, although
their magnitudes differ because the call structures differ.

The sum of the six contributions computed by leave-one-out analysis accounts
for 91.4--99.999\% of the total before-after difference observed within the same
experiment. We attribute the remaining difference to interactions among
optimizations rather than assigning it arbitrarily to one optimization.

\FloatBarrier
\subsection{Additional Techniques}

Starting from the six-optimization diagnostic used in Sect.~5.2, we enabled the
MVE cswap and paired inversion separately and together. The values in
Table~\ref{tab:additional-techniques} come from this separate diagnostic and are
not added to the final direct comparison.
\begin{table}[t]
\centering
\caption{End-to-end improvements obtained by adding MVE cswap and paired
inversion to the six-optimization diagnostic.}
\label{tab:additional-techniques}
\footnotesize
\begin{tabular}{@{}llr@{}}
\toprule
Candidate & Applicable operation & Complete-API improvement \\
\midrule
MVE cswap & Encapsulation & 1.264--1.265\% \\
Paired inversion & Encapsulation & 2.957\% \\
Both techniques & Encapsulation & \textbf{4.226--4.228\%} \\
MVE cswap & Warm decapsulation & 1.707--1.709\% \\
MVE cswap & Cold decapsulation & 0.928\% \\
\bottomrule
\end{tabular}
\end{table}

Encapsulation with both optimizations is faster than encapsulation with paired
inversion alone, so the two optimizations do not eliminate exactly the same
cost. Paired inversion applies only to encapsulation, where two outputs coexist;
key generation contains no applicable site for either optimization.

A separate candidate reordered the MVE cswap instructions to account for the
beat-level execution characteristics of MVE and reduced the function cost from
62 cycles to 54 cycles. Its end-to-end improvement for cold decapsulation was
only 0.162--0.163\%, below the predefined 0.20\% threshold, so this candidate
was excluded from the final integrated implementation.

\FloatBarrier
\subsection{Resource Cost}

We measured time and space in separate builds of the six-optimization
diagnostic after removing the experimental implementation selector and
instrumentation to approximate a deployable library. Table~\ref{tab:resource-cost}
presents the results.
\begin{table}[ht]
\centering
\caption{Time-space cost of a separate six-optimization build approximating a
deployable library. The memory increase in this table is not the exact size of
the complete final integrated implementation.}
\label{tab:resource-cost}
\footnotesize
\renewcommand{\arraystretch}{1.10}
\begin{tabularx}{\textwidth}{@{}l>{\raggedright\arraybackslash}X@{}}
\toprule
Measurement & Cost or effect \\
\midrule
Nonvolatile-memory increase & 34,308 bytes \\
Static-RAM increase & 2,040 bytes \\
Fixed-base comb precomputation table & 33,024 bytes, 96.25\% of the nonvolatile-memory increase \\
Removing the precomputation table & Approximately 152,000 additional cycles in key generation, encapsulation, and cold decapsulation \\
\bottomrule
\end{tabularx}
\end{table}

The fixed-base comb is therefore not a cost-free performance improvement; it is
a tradeoff between flash usage and execution time. Products with limited memory
can consider a configuration without the comb as a separate speed-memory
option.

\FloatBarrier
\section{Limitations}\label{sec:limitations}

\subsection{Generalizability}

The primary experiments used one EK-RA8M1 board. Reflashing the same image five
times confirmed software-level repeatability, but we did not measure a second
board without the anomalous ITCM interval, another Cortex-M85 system-on-chip,
another compiler, or another clock configuration.

The reference implementation is not an external reference point that guarantees
the highest performance among all publicly available Cortex-M85 \xwing{}
implementations. The comparison between the reference and final integrated
implementations quantifies the combined effect of the eight optimizations within
the same implementation lineage. We do not construct performance rankings from
external cycle counts obtained with different devices, APIs, or protocol
versions.

\subsection{Resource Measurement}

We measured the flash increase caused by the fixed-base table and checked for
stack-boundary corruption. An exact measurement of the code and data size of
the final integrated library still requires a build that removes all
experimental selector and instrumentation code. We did not measure power or
energy and therefore do not claim that a cycle reduction implies an equal
percentage reduction in energy.

\subsection{Security Validation}

The implementations passed KATs, byte-for-byte equivalence checks, invalid-
ciphertext rejection tests, and low-order input-point tests. In Welch's $t$-test
comparing execution times for fixed and random ciphertexts, all 24 evaluated
$t$ values had absolute values below 2.0, and the test did not detect an
execution-time difference between the two input groups. This result does not
replace a formal constant-time proof or a power or electromagnetic side-channel
evaluation.

\section{Conclusion}\label{sec:conclusion}

This work goes beyond porting \xwing{} components individually to Cortex-M85.
It jointly optimizes ML-KEM data parallelism and the fixed-base scalar
multiplication, field arithmetic, conditional swap, and inversion structure of
X25519 across the complete protocol. Direct comparison of the reference and
final integrated implementations in one executable showed cycle reductions of
19.52\% for key generation, 19.75\% for encapsulation, 11.66\% for warm
decapsulation, and 15.48\% for cold decapsulation.

The largest individual contribution was the fixed-base comb with its
33,024-byte precomputation table. MVE implementations of the ML-KEM NTT storage
order and data-format transformations, X25519 field-arithmetic rescheduling, and
common memory optimization reduced the remaining cost. The final integrated
implementation adds MVE cswap and zero-safe paired inversion to the
six-optimization configuration.

These results show that embedded cryptographic optimization requires attention
not only to faster individual functions but also to where an optimization
applies, its additional resource cost, and remeasurement at the complete API
level. Future work will remeasure end-to-end performance on other Cortex-M85
systems-on-chip, other Arm M-profile cores, and different compiler and memory
configurations to assess reproducibility and architecture dependence. It will
also measure the exact flash and RAM usage of a deployable build without
instrumentation.

\bibliographystyle{splncs04}
\bibliography{references}

\end{document}
